\section{Finite relational witnesses and backward viability}
\label{sec:viability}

A deterministic witness is often the wrong intermediate object.  At finite precision,
there may be several equally valid roots, charts, lifts, combinatorial fillings, or
approximate optimizers.  Forcing an early choice can make a correct construction
nonextensional or can smuggle in a discontinuous decision.  \(\FRNF\) therefore treats
a finite relation as a first-class output.

\subsection{Layered finite relations}

Fix a horizon \(h\).  At level \(i\le h\), candidates are numbered
\(a<B_i\).  Let \(V_i(a)\in\B\) be the validity predicate and let
\(E_i(a,b)\in\B\) express compatibility between candidate \(a\) at level \(i\)
and candidate \(b\) at level \(i+1\).  All bounds and predicates are primitive
recursive in the external parameters and final demand.

A complete path is a sequence \(a_0,\ldots,a_h\) with
\[
  V_i(a_i)=1\quad(i\le h),
  \qquad E_i(a_i,a_{i+1})=1\quad(i<h).
\]
The path need not be unique.

\subsection{Backward viability recurrence}

Define viability masks backward:
\[
\begin{split}
 W_h(a)&=V_h(a),\\
 W_i(a)&=V_i(a)\cdot
   \sg\!\left(\sum_{b<B_{i+1}}E_i(a,b)W_{i+1}(b)\right).
\end{split}
\]
Thus \(W_i(a)=1\) exactly when \(a\) can be extended to a valid path through level
\(h\).  After computing the masks, choose \(a_0\) as the least viable candidate and,
recursively, choose \(a_{i+1}\) as the least viable successor compatible with
\(a_i\).

\begin{theorem}[Backward viability]
\label{thm:backward-viability}
For every finite layered relation:
\begin{enumerate}[label=(\roman*),leftmargin=2.4em]
\item \(W_i(a)=1\) if and only if \(a\) extends to a valid path from level \(i\) to
level \(h\);
\item a complete path exists if and only if some \(a<B_0\) has \(W_0(a)=1\);
\item if a complete path exists, least bounded search applied to the masks and edges
constructs a complete path; and
\item the masks and selected path are primitive recursive and their correctness is
provable in \(\iPRA\) by quantifier-free induction.
\end{enumerate}
\end{theorem}

\begin{proof}
For (i), use backward induction on \(h-i\).  At level \(h\), viability is exactly
validity, which is the condition for the one-element suffix.  At level \(i<h\), the
recurrence says that \(a\) is valid and there exists a compatible \(b\) with
\(W_{i+1}(b)=1\).  By the induction hypothesis, that \(b\) extends to a valid suffix;
prepending \(a\) gives a valid path.  Conversely, the second element of any valid path
is a compatible viable successor, so the bounded sum has positive sign.

Clause (ii) is the case \(i=0\) of (i).  For (iii), choose the least viable initial
candidate.  If \(a_i\) is viable, the defining recurrence guarantees a compatible
viable successor; least bounded search finds one.  Induction constructs the complete
path and preserves validity and compatibility.

For (iv), encode the level-candidate table as a finite array.  The backward pass has
clock \(h+1\), and each mask entry is a bounded sum of Boolean products.  The forward
pass also has clock \(h+1\) and uses the checked least-search theorem.  The invariant
states, in Booleanized form, that processed masks satisfy the recurrence and the
selected prefix is valid, compatible, and viable.  Every quantifier is bounded by the
encoded table sizes and is therefore replaced by a characteristic equation using
\cref{thm:booleanization}.  Quantifier-free induction closes both passes.
\end{proof}

\subsection{Why no infinite choice is used}

The theorem is finite for each requested horizon.  It does not appeal to K\"onig's
lemma, dependent choice, or an infinite branch through the layered graph.  This is
crucial.  Many applications consume only a finite horizon, and stable-observer
arguments can show that the final result is independent of which finite path is used.
When an actual represented limit is required, coherence is proved quantitatively from
shrinking diameters and motion budgets, as in \cref{sec:stability}; it is not obtained
by selecting an abstract infinite branch.

\subsection{Relational composition}

Suppose a producer returns a finite relation \(R(x,a)\), while a consumer at state
\(s\) may accept several answers and move to \(s'=T(s,a)\).  A locally valid answer
can lead to a dead end at a later stage.  Backward viability solves this global finite
constraint problem: candidate \((s,a)\) is retained only when it has a compatible
continuation to the fixed final horizon.  The forward selector then produces a
successful interaction trace.

This is the correct finite replacement for a higher-type oracle strategy.  The entire
answer tree is finite after the final demand and clock are fixed.  The backward pass
computes a winning region; the forward pass selects one witness.  If the consumer's
observable is stable over all winning traces, the internal selection is erased.

\subsection{Applications}

The same recurrence covers several apparently different situations.

\begin{itemize}
\item \textbf{Root clouds.}  Candidates are rational disks or boxes with residual
certificates; edges express overlap or controlled movement.
\item \textbf{Charted geometry.}  Candidates are local charts or transition branches;
edges express overlap compatibility.
\item \textbf{Simplicial observation.}  Candidates are vertices in a carrier cloud;
edges express admissible selections.
\item \textbf{Approximate minimizers.}  Candidates are net points within a tolerance;
edges enforce motion between refinements.
\item \textbf{Finite gluing.}  Candidates are local sections; compatibility is equality
or witnessed equality on overlaps.
\item \textbf{Adaptive proof search.}  Candidates are local certificates; edges express
that the next verifier state accepts them.
\end{itemize}

The domain-specific work lies in defining validity and compatibility and proving
nonemptiness.  The viability compiler is generic.
